\documentclass[12pt]{ctexart}

\usepackage[a4paper,margin=2.54cm]{geometry}
\usepackage{amsmath,amssymb}
\usepackage{booktabs}
\usepackage{threeparttable}
\usepackage{array}
\usepackage{graphicx}
\usepackage{hyperref}

\title{题目待定}
\author{作者待填}
\date{\today}

\begin{document}

\maketitle

\begin{abstract}
摘要待补。建议按“研究问题、数据与方法、主要发现、贡献与启示”的顺序填写。
\end{abstract}

\noindent\textbf{关键词：} 关键词1；关键词2；关键词3

\section{引言}
引言待补。建议按“现实背景、研究缺口、研究问题、研究设计、研究贡献、文章结构”的顺序展开。

\section{文献综述与研究假设}
文献综述待补。建议按主题或机制组织，而不是按作者逐篇罗列。

\section{数据来源、变量说明与模型设定}
数据来源、样本处理、变量构造和模型设定待补。

\section{实证结果}
结果概述待补。
% 主结果表接口示例
% \input{../05_analysis_output/export_for_paper/paper_table_01_descriptive_statistics.tex}
% \input{../05_analysis_output/export_for_paper/paper_table_03_baseline_model.tex}

\section{稳健性、异质性或机制分析}
补充分析待补。

\section{讨论}
讨论待补。重点解释主要结果、反直觉结果和研究边界。

\section{结论与启示}
结论待补。重点回到研究问题，总结发现，说明启示与局限。

% \bibliographystyle{plain}
% \bibliography{references}

\end{document}
